\section{Introduction}
Large language models (LLMs) are increasingly deployed in high-stakes settings where users may challenge, pressure, or contradict the model's answers. A reliable model should revise its stance when confronted with genuine evidence, but should not change its stance simply because a user expresses displeasure or repeats an assertion with greater confidence. Sycophancy is the alignment of a model with the user's stated belief or preference even when doing so conflicts with truthfulness~\citep{perez2022discoveringlanguagemodelbehaviors,sharma2023towardsunderstanding,ranaldi2023whenlargelanguagemodels}. Prior work shows that this behavior is not incidental: models trained with human feedback can learn to favor responses that match user views, and both human raters and preference models may reward convincing agreement over correctness~\citep{sharma2023towardsunderstanding,wei2023simplesyntheticdata}. Sycophancy can be exhibited even in chain-of-thought reasoning, where traces can be biased by social cues the user introduces, with the model's stated reasoning shifting to justify the cued answer rather than the correct one~\citep{turpin2023languagemodelsdontalways}.

Recent work has established sycophancy as an interaction-level failure where multi-turn benchmarks such as TRUTH DECAY~\citep{liu2025truthdecayquantifyingmultiturn} and SYCON Bench~\citep{hong2025syconbench} show that models can drift or flip under sustained user pressure, while SWAY~\citep{bhalla2026sway} shows that agreement can shift under counterfactual linguistic framing. These works demonstrate that sycophancy in multi-turn evaluation reveals failures missed by single-turn tests. Existing benchmarks explore sycophancy through the flip in stance and its frequency.

Motivated by these works, we ask whether a model's uncertainty on a specific question predicts how it responds to escalating pressure from the user, and whether this response is visible in the model's reasoning trajectory. Our study is organized around four research questions:
\begin{enumerate}[leftmargin=*, labelwidth=3em, labelsep=0.5em, align=left]
    \item[\textbf{RQ1:}] Does pre-pressure uncertainty moderate susceptibility to sycophancy?
    \item[\textbf{RQ2:}] Does pressure dose, conditioned on uncertainty, predict sycophantic flipping?
    \item[\textbf{RQ3:}] Do uncertainty-conditioned pressure effects vary across models?
    \item[\textbf{RQ4:}] Do chain-of-thought (CoT) belief trajectories reflect capitulation?
\end{enumerate}

To answer these research questions, we develop a multi-turn evaluation framework that estimates model's uncertainty through repeated sampling, then measures answer flips under escalating pressure. We evaluate five OpenAI and Anthropic models on MMLU-Pro~\citep{wang2024mmluprorobustchallengingmultitask} and GPQA~\citep{rein2023gpqagraduatelevelgoogleproofqa}, and add Gemini 3.5 Flash on Humanity's Last Exam (HLE)~\citep{hle} and the single-turn category comparison to understand the effect of uncertainty and hardness on the eventual flip of stance from the model and if the eventual flip could be identified through the reasoning trace of these models.